% === I. INTRODUCTION =========================================================
\section{Introduction}

\IEEEPARstart{M}{odel-based} controllers for autonomous racing derive their
performance from the fidelity of the vehicle model they carry
\cite{liniger2015,kabzan2019lbmpc,forzaeth2024}. In the lateral channel that
fidelity is dominated by the tire, simultaneously the physical bound on
manoeuvrability and the least accessible component to measure: its behaviour
depends on surface, temperature, wear and load in ways that change between
sessions \cite{pacejka2012book,amz2020}.

Classical characterisation resolves this by taking the vehicle off the track.
Steady-state circular or ramp-steer manoeuvres in a large open space yield
clean data, at a logistical cost a race weekend rarely affords
\cite{ontrack2025}. Identifying on track removes the logistics and brings back
transients, noise and a racing line that is a poor excitation signal for a
nonlinear model; nonlinear least squares applied directly to lap data is
brittle in exactly this regime \cite{ontrack2025}.

Dikici \emph{et al.}~\cite{ontrack2025} propose the hybrid remedy that is the
starting point of this paper. A small residual network is trained on
\SI{30}{\second} of ordinary lap data to predict the one-step error of a
nominal single-track model. The corrected model is rolled out through a
\emph{synthetic}, quasi-steady steering sweep, and Pacejka coefficients are
fitted to that sweep. Because the sweep is generated rather than driven it can
be made steady and noise-free, and because the residual network only has to
interpolate within its training distribution, its generalisation weakness stays
contained. The procedure iterates, each fit becoming the next nominal model. On
a 1:10 platform it reports a \num{3.3}$\times$ lower one-step RMSE than
nonlinear least squares under noise, from \SI{30}{\second} of driving
\cite{ontrack2025}.

\subsection{Problem statement}

This paper asks what the pipeline identifies when it is moved from a
\SI{3.7}{\kilo\gram}, \SI{0.32}{\meter}-wheelbase platform to a full-scale
Formula Student vehicle: \SI{269.6}{\kilo\gram}, \SI{1.53}{\meter} wheelbase,
\SI{16}{\degree} of road-wheel lock, following a reference speed profile of
\SIrange{9.6}{25.8}{\meter\per\second} and achieving
\SIrange{9.2}{26.5}{\meter\per\second} in the identification data. The
question is identifiability, not parameter transfer, and two scale-dependent
mechanisms, both latent at 1:10, break the estimate:

\begin{enumerate}
\item \textbf{Excitation collapses inside the virtual-data step.} The Pacejka
fit never sees the recorded lap. It sees the synthetic rollout, whose reachable
friction $\mu_{\mathrm{reach}}\!\approx\!v^{2}\delta_{\mathrm{sweep}}/(gL)$ is a
property of the rollout configuration and not of the tire, and a rollout speed
inherited from the scaled platform caps it at $0.43$ here. Below the tire's
peak the Magic Formula degenerates to $F_{y}\approx F_{z}(BCD)\alpha$, only
the product $BCD$ is identifiable, and the bounded optimiser resolves the
remaining freedom on a box edge \cite{ljung1999sysid,pacejka2012book}.

\item \textbf{Lateral demand scales as $v^{2}$, so a bad grip estimate is
silent until it is catastrophic.} Driven through the same controller on an
earlier, faster reference trajectory, one identified set with a peak friction
coefficient of \num{0.40} tracked at \SI{0.16}{\meter} of cross-track error at
\SI{11.1}{\meter\per\second}, \SI{0.43}{\meter} at
\SI{15}{\meter\per\second} and \SI{41.89}{\meter} at
\SI{27}{\meter\per\second}. The scaled platform never reaches the speed at
which the error becomes observable. These figures come from an offline
closed-loop harness rather than a controlled simulator run and are reported as
such throughout (Section~\ref{subsec:gates}).
\end{enumerate}

One difference works in our favour. Because the study runs in the CARLA
simulator \cite{carla2017} through a purpose-built ROS~2 interface, the
simulator's own per-wheel telemetry (slip angle, normal load and lateral force)
is available as ground truth, so an on-track identifier can be scored against
the forces the plant actually generated.

\subsection{Contributions}

\begin{enumerate}
\item \textbf{An identifiability analysis of the virtual-data step}
(Section~\ref{subsec:excitation}), giving the closed-form reachable-friction
bound above and an \emph{excitation-aware rollout} that selects the rollout
speed from the measured lap and warns whenever the configuration cannot reach
the prior's peak. A grid over rollout speed, sweep amplitude and true peak
factor on a known tire verifies it (Section~\ref{subsec:sweep}): below the
tire's own peak the fit returns the excitation rather than the tire, at every
grip level tested and independently of the solver. This corrects the baseline
method~\cite{ontrack2025} at any scale, and stays latent at 1:10 because a
\SIrange{2}{4}{\meter\per\second} rollout is representative there and is not
here.

\item \textbf{Two structural properties of the loop, and the corrections they
imply} (Section~\ref{subsec:contraction}). Because stage~3 reconstructs its
forces from the rollout's own quasi-steady yaw balance, it observes nothing but
the lateral acceleration the previous coefficient set produced, so its residual
depends on $B$, $C$ and $D$ only through their product: the peak factor is not
recoverable from the rollout \emph{at any excitation}, a stronger statement
than the low-slip degeneracy, and must instead be measured on the recorded
trace and held. And because each fit is handed the trajectory its own previous
answer generated, the loop is a positive feedback path rather than a
contraction, and iterating it undamped degrades the estimate monotonically.
Stage~4 becomes a relaxed fixed-point iteration at $\beta=\num{0.20}$; the
$\beta=0$ control, which performs no identification at all, is the worst
setting of the sweep, so the iterations do carry information.

\item \textbf{The corrections the bound implies}, on both sides of the fit.
On the rollout side (Section~\ref{subsec:iterative}), the sweep is bounded at
$1.5\times$ the 99th percentile of the observed steering, so the fit cannot
read the residual network's extrapolation as grip, and the Tikhonov target is
frozen at the configured prior while the optimiser's start point tracks the
iterations; together those remove a drift that took $D$ from the
\num{1.009}/\num{1.002} front/rear prior to \num{1.5402} and \num{1.8721} over
six iterations. On the measurement side (Section~\ref{subsec:measurement}),
slip angles are built from the \emph{achieved} road-wheel angle rather than
the command, mass and wheelbase from measured static wheel loads rather than
from a configuration file, the friction warm start from the \emph{curvature} of
the axle force curve rather than a median over a low-slip mask, and the
dynamics residuals of the physics-informed objective are normalised to the
state-increment units of the data term, which the published form leaves a
factor \num{2.3e5} apart.

\item \textbf{Admissibility gates on derived physical quantities}
(Section~\ref{sec:integration}) at the model-to-controller boundary: axle
cornering stiffness, understeer sign and critical speed, and grip ceiling
against the trajectory's lateral demand, rather than per-coefficient box
checks, which cannot detect a degenerate fit or a pairwise-oversteering axle
set.

\item \textbf{A simulation architecture that supplies the ground truth, and a
closed-loop evaluation on it} (Sections~\ref{subsec:platform}--\ref{subsec:data}
and \ref{sec:results}). \bridge{} exposes the plant's own per-wheel forces,
slip angles and loads as ROS~2 topics alongside a conventional autonomy
interface. The evaluation runs on a static \num{0.70} friction surface,
deliberately not the \num{1.05} the configured tire prior was measured on, so
that a pipeline unable to move off its prior has nowhere to hide
(Section~\ref{subsubsec:priorprovenance}). Over three timed laps the corrected
pipeline identifies axle cornering stiffness to \SI{3.2}{\percent} front and
\SI{1.4}{\percent} rear and tracks the peak factor to \SI{11}{\percent},
against an inherited configuration that rails $D$ on its \num{2.0} bound in
every cycle and overstates stiffness by \SIrange{118}{132}{\percent}. Axle
force RMSE against the plant's telemetry is \num{12.1}$\times$ and
\num{10.5}$\times$ lower, and MPC lateral tracking error \SI{32}{\percent}
lower, at unchanged solver cost.
\end{enumerate}

All results are obtained in simulation, and Section~\ref{sec:disc} states which
conclusions are expected to transfer to hardware.
